% 322_Spektraltheorie_Hilbert_En_ch.tex
% Spectral Theory and Hilbert-Space Embedding of the FFGFT Tori
% Mathematical Foundation of the T^4/Z_3 Spectral Mechanics
% Doc. 322, August 2026

% =========================================================
\section*{List of Symbols}
% =========================================================

\begin{center}
\begin{tabular}{lll}
\toprule
Symbol & Meaning & Value / Unit \\
\midrule
$\mathcal{H}$ & Hilbert space of physical states & separable \\
$\hat{F}$ & fundamental fractal spectral operator & self-adjoint \\
$\hat{F}_{D_4}$ & $D_4$ sub-operator & eigenvalue $\xi$ \\
$\sigma(\hat{F})$ & spectrum of $\hat{F}$ & $\subset \mathbb{R}$ \\
$E(\lambda)$ & spectral projection measure & $E(\lambda)^2 = E(\lambda)$ \\
$\mathcal{D}(\hat{F})$ & domain of $\hat{F}$ & dense in $\mathcal{H}$ \\
$P_{\mathbb{Z}_3}$ & $\mathbb{Z}_3$ projection operator & idempotent \\
$\hat{\Delta}_{T^4}$ & Laplacian on $T^4/\mathbb{Z}_3$ & essentially self-adj. \\
$\hat{W}$ & winding-number operator & eigenvalues $(n_\theta, n_\phi)$ \\
$D_f$ & fractal dimension of the orbifold & $3 - \xi$ \\
$D_f^{\mathrm{frac}}$ & approximate fractal dimension & $2.94$ (no $\xi$-derivation) \\
$K_{\mathrm{frac}}$ & interface factor rolled $\to$ SI & $1-100\xi = 74/75$ (Doc.~133) \\
$\xi$ & FFGFT geometry parameter & $4/30000$ \\
$\omega$ & $\mathbb{Z}_3$ phase factor & $e^{2\pi i/3}$ \\
$[K]$ & derivable from $\xi$ & status marker \\
$[B]$ & algebraically proved & status marker \\
$[S]$ & sketched / plausible & status marker \\
$[\text{SETZUNG}]$ & axiomatic setting & status marker \\
\bottomrule
\end{tabular}
\end{center}

% =========================================================
\section*{Summary}
% =========================================================

Doc.~322 develops the mathematical foundation of FFGFT spectral theory.
The mass formulae of Doc.~320 receive here their rigorous Hilbert-space
embedding: the fundamental fractal operator $\hat{F}$ on the state space
$\mathcal{H}_{\mathrm{FFGFT}}$ has the particle masses as eigenvalues.
The geometry parameter $\xi = 4/30000$ is the smallest eigenvalue of the
$D_4$ sub-operator. The mass formulae of Docs.~006, 317 and 320 are not
phenomenological ans\"{a}tze but eigenvalue equations of a self-adjoint
operator.

\begin{keyresult}[Status markers]
All statements carry FFGFT status markers:\\
\textbf{[K]} derivable from $\xi$ \quad
\textbf{[B]} algebraically proved \quad
\textbf{[S]} sketched/plausible \quad
\textbf{[SETZUNG]} axiomatic setting
\end{keyresult}

% =========================================================
\section{Foundations of Spectral Theory in Hilbert Spaces}
% =========================================================

\subsection{Motivation: from matrices to operators}

In linear algebra a symmetric $n\times n$ matrix $A$ possesses exactly
$n$ real eigenvalues and an orthonormal eigenbasis, enabling the
diagonal decomposition $A = UDU^T$.

In quantum mechanics -- and in FFGFT -- operators act on
\emph{infinite-dimensional} Hilbert spaces. The momentum operator,
the Hamiltonian, and the FFGFT spectral operator have no
finite-dimensional analogue. Spectral theory generalises the
eigenvalue concept to this case.

\subsection{The Hilbert space: definition and physical meaning}

\textbf{Definition (Hilbert space).}
A Hilbert space $\mathcal{H}$ is a complete inner-product vector space.
Its elements have finite norm $\|\psi\| = \sqrt{\langle\psi,\psi\rangle}
< \infty$, which in position space means
\begin{equation}
  \int_{-\infty}^{\infty} |\psi(x)|^2 \, dx < \infty.
\end{equation}
Only such square-integrable wave functions are physically realisable --
they describe normalisable probability densities.

\textbf{Note:} Eigenstates of operators with \emph{continuous} spectrum
-- e.g.\ position eigenstates $|y\rangle$ with $\langle x|y\rangle =
\delta(x-y)$ -- are not normalisable and lie outside $\mathcal{H}$
itself, in the extended distributional space. In FFGFT this issue
arises in the IR limit; the sub-Planck regime has a purely discrete
spectrum.

\subsection{The spectrum of a linear operator}

For a linear operator $T$ on a Hilbert space the spectrum $\sigma(T)$
generalises the set of all eigenvalues:
\begin{equation}
  \sigma(T) = \{\lambda \in \mathbb{C} : (T - \lambda I)
  \text{ not invertible with bounded inverse}\}.
\end{equation}
One distinguishes:
\begin{itemize}
  \item \textbf{Point spectrum $\sigma_p(T)$:} ``true'' eigenvalues
    with $T\psi = \lambda\psi$, $\psi \neq 0$.
  \item \textbf{Continuous spectrum $\sigma_c(T)$:}
    $(T-\lambda I)$ injective but not surjective; no normalised
    eigenvector.
  \item \textbf{Absolutely continuous and singular spectrum:}
    important for Schr\"{o}dinger operators on fractal spaces.
\end{itemize}

\subsection{The spectral theorem: centrepiece of the theory}

\textbf{[B]} For self-adjoint operators ($T = T^\dagger$ with suitable
domain) the spectral theorem holds: there exists a unique
\emph{spectral measure} $E$ -- a projection-valued function
$\lambda \mapsto E(\lambda)$ -- such that
\begin{equation}
  T = \int_{\sigma(T)} \lambda \, dE(\lambda).
  \label{eq:spectralthm}
\end{equation}
For compact operators Eq.~\eqref{eq:spectralthm} reduces to the
familiar discrete sum $T = \sum_i \lambda_i \langle\cdot, e_i\rangle e_i$.
Self-adjointness guarantees $\sigma(T) \subset \mathbb{R}$ -- real
spectra, the necessary condition for measurable observable values.

% =========================================================
\section{The Fundamental Fractal Operator $\hat{F}$}
% =========================================================

\subsection{Definition and Hilbert-space structure}

\textbf{[SETZUNG]} FFGFT postulates a quantum-logical fractal operator
$\hat{F}$ acting on a separable Hilbert space $\mathcal{H}$, formally
encoding the fractal self-similarity of the sub-Planck geometry:
\begin{equation}
  \hat{F} \equiv \bigoplus_{n=1}^{\infty} S_n \circ L_n,
  \label{eq:F_def}
\end{equation}
with scale operators $S_n = r_n U_n$ ($r_n \in (0,1)$ contraction
rates, $U_n$ unitary) and logical projectors $L_n$ satisfying
$L_n^2 = L_n$, $L_n^\dagger = L_n$, $L_nL_m = L_{\min(n,m)}$.

The domain
\begin{equation}
  \mathcal{D}(\hat{F}) = \left\{ \psi \in \mathcal{H} :
  \sum_{n=1}^{\infty} \|S_n(L_n\psi)\|^2 < \infty \right\}
\end{equation}
is dense in $\mathcal{H}$. \textbf{[B]} $\hat{F}$ is self-adjoint:
since $L_0 = \xi\cdot\ell_P$ (Doc.~180) renders the scale ladder finite,
$\sum_n r_n < \infty$ and hence $\hat{F}$ is bounded with
$\mathcal{D}(\hat{F}) = \mathcal{H}$; symmetry follows from the
$\mathbb{Z}_3$ pairing $(k,-k)$. The deficiency indices are $(0,0)$ --
there is exactly one self-adjoint realisation (Doc.~327, R86).

\subsection{$\xi$ as smallest eigenvalue of the $D_4$ sub-operator}

\textbf{[K]} The geometry parameter is not a free parameter but a
spectral value:
\begin{equation}
  \xi = \lambda_{\min}(\hat{F}_{D_4}) = \frac{4}{30000}.
  \label{eq:xi_eigen}
\end{equation}
This explains structurally why $\xi$ appears as the base of all mass
formulae: it is not a chosen parameter but the lowest mode of the
geometry operator.

% =========================================================
\section{The FFGFT State Space}
% =========================================================

\subsection{Tensor-product structure}

\textbf{[K]} The physical state space:
\begin{equation}
  \mathcal{H}_{\mathrm{FFGFT}}
  = \mathcal{H}_{\mathrm{geom}}
  \otimes \mathcal{H}_{\mathrm{spin}}
  \otimes \mathcal{H}_{\mathrm{flavor}},
\end{equation}
with:
\begin{enumerate}
  \item $\mathcal{H}_{\mathrm{geom}} = L^2(T^4/\mathbb{Z}_3, d\mu_f)$:
    square-integrable functions on the fractal orbifold.
  \item $\mathcal{H}_{\mathrm{spin}} = \mathbb{C}^{2\lceil D_f\rceil}$:
    spinor space; $\lceil D_f \rceil = 3$.
  \item $\mathcal{H}_{\mathrm{flavor}} = \ell^2(\mathbb{N}^2)$:
    flavour space with winding numbers $(n_\theta, n_\phi)$.
\end{enumerate}

\subsection{The fractal measure $d\mu_f$ with 100-fold recursion}

\textbf{[K]} The fractal measure on $T^4/\mathbb{Z}_3$ is built by
100 successive fractalisations of the Lebesgue measure:
\begin{align}
  d\mu_f^{(0)}(x) &= dx, \\
  d\mu_f^{(k+1)}(x) &= \prod_{i=1}^4
  \!\left(\sum_{n=1}^{100} c_n^{(k)}\,
  \delta\!\left(x_i - x_i^{(n)}\right)\right)
  d\mu_f^{(k)}(x),
\end{align}
with $c_n^{(k)} = n^{-(D_f^{(k)}-3)}$,
$D_f^{(k)} = 3-\xi^{(k)}$,
$\xi^{(k)} = (4/30000)(1-k/100)$.
After 100 steps:
\begin{equation}
  \boxed{d\mu_f(x)
  = \prod_{i=1}^4\!\left(\sum_{n=1}^{100} n^{\xi}\,
  \delta\!\left(x_i-x_i^{(n)}\right)\right)dx_i},
  \quad \xi = \tfrac{4}{30000}.
\end{equation}

% =========================================================
\section{The Torus $T^4$ in Hilbert Space}
% =========================================================

\subsection{Fourier decomposition on $T^4$}

\textbf{[B]}
\begin{equation}
  L^2(T^4) \cong \bigoplus_{n\in\mathbb{Z}^4}
  \mathbb{C}\cdot e^{in\cdot x/R},\qquad
  \Delta_{T^4}\,e^{in\cdot x/R} = -\tfrac{|n|^2}{R^2}\,e^{in\cdot x/R}.
\end{equation}

\subsection{The $\mathbb{Z}_3$ orbifold projection}

\textbf{[K]}
\begin{equation}
  P_{\mathbb{Z}_3} = \tfrac{1}{3}\sum_{k=0}^2 \omega^k\tau^k,
  \qquad
  \mathcal{H}(T^4/\mathbb{Z}_3) = \mathrm{Fix}(P_{\mathbb{Z}_3})
  \subset L^2(T^4).
\end{equation}

\subsection{Fixed points and neutrino localisation}

\textbf{[K]} The boundary condition at a fixed point $x_0$:
\begin{equation}
  \psi_\chi(x_0 + y) = \chi\cdot\psi_\chi(x_0 + g_*(y)).
  \label{eq:rbd}
\end{equation}
The three neutrino flavours are localised at fixed points
(Doc.~320):
\begin{equation}
  \nu_1 \leftrightarrow F_2 = \tfrac{2\pi R}{3}(1,1,1,0),\quad
  \nu_2 \leftrightarrow F_5 = \tfrac{4\pi R}{3}(1,1,1,0),\quad
  \nu_3 \leftrightarrow F_7 = \tfrac{4\pi R}{3}(1,1,1,1).
\end{equation}

% =========================================================
\section{The Complete Spectral Operator}
% =========================================================

\subsection{Tensor-product form}

\textbf{[K]}
\begin{equation}
  \hat{F} = \hat{\Delta}_{T^4} \otimes \hat{S} \otimes \hat{W},
  \qquad
  \hat{M}_f = v\cdot\xi^{\hat{p}} \otimes \hat{r}.
\end{equation}

\subsection{Commutator conditions and MASA}

\textbf{[K]}
\begin{equation}
  [\hat{F}, \hat{S}] = 0,\qquad
  [\hat{n}_\theta, \hat{n}_\phi] = 0.
\end{equation}
The MASA $\mathcal{A} = \mathrm{span}\{\hat{n}_\theta, \hat{n}_\phi,
\hat{p}\}$ yields the complete spectral labelling $(n_\theta, n_\phi, p)$.

% =========================================================
\section{Generalised Fourier Transform on $T^4/\mathbb{Z}_3$}
% =========================================================

\textbf{[K]} Projection kernel and GFT:
\begin{equation}
  K_\lambda^{(p)}(x) = \sum_{n\in\mathbb{Z}^4}
  \frac{\delta_{\lambda,\Lambda_n^{(p)}}}{|n|^{D_f/2}}\,e^{in\cdot x/R},
  \qquad
  \Lambda_n^{(p)} = \frac{|n|^2}{R^2}\cdot\xi^p.
\end{equation}
\begin{equation}
  \tilde{\psi}(\lambda) = \langle K_\lambda^{(p)}, \psi\rangle_{\mathcal{H}}
  = \int_{T^4/\mathbb{Z}_3}
  \overline{K_\lambda^{(p)}(x)}\,\psi(x)\,d\mu_f(x).
\end{equation}

% =========================================================
\section{Spectral Derivation of Particle Masses}
\label{sec:masses}
% =========================================================

\subsection{Charged leptons as torus excitations}

\textbf{[K]} The mass eigenvalue equation
$\hat{M}_{e,\mu,\tau}|n_\theta,n_\phi,p\rangle = m_i|n_\theta,n_\phi,p\rangle$
gives:
\begin{equation}
  m_i = \frac{n_\phi^2}{n_\theta}\cdot\xi^{p_i}\cdot v.
\end{equation}
\begin{center}
\begin{tabular}{lcccc}
\toprule
Particle & $n_\theta$ & $n_\phi$ & $p_i$ & $m_i$ \\
\midrule
Electron $e$  & 3 & 2 & $3/2$ & $0.511$~MeV \\
Muon $\mu$    & 5 & 4 & $1$   & $104.96$~MeV \\
Tau $\tau$    & 9 & 5 & $2/3$ & $1783.5$~MeV \\
\bottomrule
\end{tabular}
\end{center}

\subsection{Neutrinos as fixed-point eigenstates}

\textbf{[K]}
\begin{equation}
  P_{\mathbb{Z}_3}|\nu_i\rangle = |\nu_i\rangle,\qquad
  m_{\nu_i} = m_e\cdot\xi^{p_i}.
\end{equation}
\begin{center}
\begin{tabular}{lcccc}
\toprule
Neutrino & Fixed point & $\dim_{\mathrm{loc}}$ & $p_i$ & $m_{\nu_i}$ \\
\midrule
$\nu_1$ & $F_2$ & $2-1/4$ & $9/4$ & $0.976$~meV \\
$\nu_2$ & $F_5$ & $2$     & $2$   & $9.084$~meV \\
$\nu_3$ & $F_7$ & $2-1/5$ & $9/5$ & $54.11$~meV \\
\bottomrule
\end{tabular}
\end{center}

\subsection{Spectral selection rule for the exponent ladder}

\textbf{[K]}
\begin{equation}
  \Delta p = \frac{k}{3},\quad k\in\mathbb{Z},\qquad
  \hat{F}\cdot\hat{\xi}^{1/3} = \hat{\xi}^{1/3}\cdot\hat{F}.
\end{equation}

% =========================================================
\section{Fractal Dimension and $K_{\mathrm{frac}}$}
% =========================================================

\textbf{[K]} Exact FFGFT formula:
\begin{equation}
  \boxed{D_f = 3 - \xi = 3 - \frac{4}{30000}
  = \frac{22499}{7500} \approx 2.9998\overline{6}}.
\end{equation}
Doc.~133 gives the interface factor:
\begin{equation}
  K_{\mathrm{frac}} = 1 - 100\xi = \frac{74}{75} \approx 0.9866\overline{6}.
\end{equation}
In mass ratios $K_{\mathrm{frac}}$ cancels exactly (Doc.~012).

% =========================================================
\section{Complete Spectral Correspondence Table}
% =========================================================

\begin{table}[h!]
\centering
\caption{All spectral states of the FFGFT and their Hilbert-space
  components (Doc.~322)}
\label{tab:spectral}
\begin{tabular}{llllll}
\toprule
Sector & State & Eigenvalue & Formula &
  $\mathcal{H}$-component & Status \\
\midrule
Leptons
  & $|e\rangle$    & $0.511$~MeV  & $\frac{4}{3}\xi^{3/2}v$
    & $\mathcal{H}_{\mathrm{fl}}(3,2)$ & [K] \\
  & $|\mu\rangle$  & $104.96$~MeV & $\frac{16}{5}\xi^1 v$
    & $\mathcal{H}_{\mathrm{fl}}(5,4)$ & [K] \\
  & $|\tau\rangle$ & $1783.5$~MeV & $\frac{25}{9}\xi^{2/3}v$
    & $\mathcal{H}_{\mathrm{fl}}(9,5)$ & [K] \\
\midrule
Neutrinos
  & $|\nu_1\rangle$ & $0.976$~meV & $m_e\xi^{9/4}$
    & $\mathrm{Fix}(P)\cap F_2$ & [K] \\
  & $|\nu_2\rangle$ & $9.084$~meV & $m_e\xi^2$
    & $\mathrm{Fix}(P)\cap F_5$ & [K] \\
  & $|\nu_3\rangle$ & $54.11$~meV & $m_e\xi^{9/5}$
    & $\mathrm{Fix}(P)\cap F_7$ & [K] \\
\midrule
Bosons
  & $|g_a\rangle$   & $0$          & $\hat{F}|g_a\rangle=0$
    & $\mathrm{Fix}(P_{\mathbb{Z}_3})$ & [K] \\
  & $|W^\pm\rangle$ & $80.38$~GeV  & $\frac{1}{2}g_2 v$
    & $\mathcal{H}_{\mathrm{spin}}$ & [B] \\
  & $|Z^0\rangle$   & $91.19$~GeV  & $m_W/\!\cos\theta_W$
    & $\mathcal{H}_{\mathrm{spin}}$ & [B] \\
\bottomrule
\end{tabular}
\end{table}

\begin{keyresult}[Overall status Doc.~322]
The Hilbert-space embedding of FFGFT is structurally complete:
$\hat{F}$ is defined, the state space is constructed, all mass
formulae from Doc.~320 are identified as eigenvalue equations, and the
self-adjointness of $\hat{F}$ is proved (Doc.~327, R86):
$\xi = \lambda_{\min}(\hat{F}_{D_4})$ is thereby well-defined as a
spectral value.
The particle spectrum follows from spectral theory, not from
free parameters.
\end{keyresult}
